
\subsection{\texorpdfstring{$\mathcal M_{k+1}^{\mathrm{main}}
(\{J_\rho\}_\rho:\beta;\operatorname{twp}(\rho);P_1,\ldots,P_k)$
的定向}{逐时积模空间的定向}}
\label{subsec:8-9-2}

本小节说明
$S([0,1]\times L;\Lambda_{0,\mathrm{nov}})$ 上
$\mathfrak m_{1,\beta_0}$ 的符号，以及逐时积模空间的定向。
（见第~\ref{subsec:7-2-14} 小节。）

\begin{definition}{8.9.5}\label{def:8-9-5}
对 $P_j\in S([0,1]\times L;\mathbb Q)$，定义
\translatornote{源书在下式右端把已经属于
$S([0,1]\times L;\mathbb Q)$ 的每个 $P_i$ 又写成
$[0,1]\times P_i$；这会使该因子不再自然映入纤维积底空间
$([0,1]\times L)^k$。式 (7.2.362) 对圆柱中的一般单形直接使用
$P_i$，而本式的 $n+1-\dim P_j$ 也正是 $n+1$ 维圆柱中的次数，
故删去多余的区间因子。}
\[
\begin{aligned}
 &\mathcal M_{k+1}^{\mathrm{main}}
 (\{J_\rho\}_\rho:\beta;\operatorname{twp}(\rho);
 P_1,\ldots,P_k)\\
 &\quad={}
 (-1)^\delta
 \mathcal M_{k+1}^{\mathrm{main}}
 (\{J_\rho\}_\rho:\beta;\operatorname{twp}(\rho))
 \mathop{\times}_{([0,1]\times L)^k}
 \left(\prod_{i=1}^{k}P_i\right),
\end{aligned}
\]
其中
\[
 \delta=(n+2)
 \sum_{i=1}^{k-1}\sum_{j=1}^{i}(n+1-\dim P_j).
\]
\end{definition}

逐时积模空间的定义见式~\eqref{eq:7-2-359} 与
\eqref{eq:7-2-362}。定义~\ref{def:8-9-5}
与定义~\ref{def:8-4-1} 相容。回忆，由定义~\ref{def:8-4-1}，
\[
 \mathcal M_{k+1}(\beta;P_1,\ldots,P_k)
 =(-1)^{\epsilon_1}\mathcal M_{k+1}(\beta)
 \mathop{\times}_{L^k}
 \left(\prod_{i=1}^{k}P_i\right),
\]
其中
\[
 \epsilon_1=(n+1)
 \sum_{i=1}^{k-1}\sum_{j=1}^{i}\deg P_j,
 \qquad
 \deg P_j=\dim L-\dim P_j=n-\dim P_j.
\]

本小节余下部分比较
\[
 [0,1]\times
 \mathcal M_{k+1}(J_\rho;\beta;P_1,\ldots,P_k)
\]
与
\[
 \mathcal M_{k+1}
 (\{J_\rho\}_\rho:\beta;\operatorname{twp}(\rho);
 [0,1]\times P_1,\ldots,[0,1]\times P_k)
\]
的定向。由于这里只关心符号亦即定向，以下假设 $J_\rho\equiv J$。
我们来证明下述结果。

\begin{lemma}{8.9.6}\label{lem:8-9-6}
\[
\begin{aligned}
 &\left([0,1]\times\mathcal M_{k+1}(\beta)\right)
 \mathop{\times}_{([0,1]\times L)^k}
 \left(\prod_{i=1}^{k}([0,1]\times P_i)\right)\\
 &\qquad={}
 (-1)^{\epsilon_2}[0,1]\times
 \left(\mathcal M_{k+1}(\beta)
 \mathop{\times}_{L^k}\prod_{i=1}^{k}P_i\right),
\end{aligned}
\]
其中
\[
 \epsilon_2=\sum_{i=1}^{k-1}\sum_{j=1}^{i}(n-\dim P_j).
\]
\end{lemma}

\begin{proof}
置
\[
\begin{aligned}
 X&=\mathcal M_{k+1}(\beta),
 &I&=I_0=I_1=\cdots=I_k=[0,1],\\
 L_0&=L_1=\cdots=L_k=L.
\end{aligned}
\]
于是所比较的两个空间分别写成
\[
\begin{aligned}
 &(I\times X)
 \mathop{\times}_{\prod_{i=1}^{k}(I_i\times L_i)}
 \left(\prod_{i=1}^{k}(I_i\times P_i)\right),\\
 &\left(I\mathop{\times}_{\prod_{i=1}^{k}I_i}
 \prod_{i=1}^{k}I_i\right)
 \times
 \left(X\mathop{\times}_{\prod_{i=1}^{k}L_i}
 \prod_{i=1}^{k}P_i\right).
\end{aligned}
\]
由下列等式定义
${}^\circ I^\circ$、${}^\circ P_i$、${}^\circ X^\circ$
和 ${}^\circ(I\times X)^\circ$
\translatornote{源书在最后一式的底空间因子中写
$I_i\times P_i$；但 $I\times X$ 的靶是
$(I_0\times L_0)\times\prod_i(I_i\times L_i)$，
下一页的展开也始终使用 $I_i\times L_i$。故将该处 $P_i$
恢复为 $L_i$。}：
\[
\begin{aligned}
 I&=I_0\times{}^\circ I^\circ\times\prod_{i=1}^{k}I_i,
 &P_i&=L_i\times{}^\circ P_i,\\
 X&=L_0\times{}^\circ X^\circ\times\prod_{i=1}^{k}L_i,
 &I\times X&=(I_0\times L_0)
 \times{}^\circ(I\times X)^\circ
 \times\prod_{i=1}^{k}(I_i\times L_i).
\end{aligned}
\]
注意虚维数为
\[
 \dim{}^\circ X^\circ=\dim X-(k+1)n,
 \qquad
 \dim{}^\circ I^\circ=-k.
\]
由于
\[
\begin{aligned}
 I\times X
 &=\left(I_0\times{}^\circ I^\circ\times\prod_{i=1}^{k}I_i\right)
 \times
 \left(L_0\times{}^\circ X^\circ\times\prod_{i=1}^{k}L_i\right)\\
 &=(-1)^a(I_0\times L_0)
 \times{}^\circ I^\circ\times{}^\circ X^\circ
 \times\prod_{i=1}^{k}(I_i\times L_i),
\end{aligned}
\]
其中
\[
 a=k\dim X+\frac{k(k-1)n}{2},
\]
故
\[
 {}^\circ(I\times X)^\circ
 =(-1)^a{}^\circ I^\circ\times{}^\circ X^\circ.
\]
另一方面，写成
\[
 I_i\times P_i=(I_i\times L_i)
 \times{}^\circ(I_i\times P_i).
\]
由于
$I_i\times P_i=I_i\times(L_i\times{}^\circ P_i)$，有
\[
 {}^\circ(I_i\times P_i)={}^\circ P_i.
\]
因此得到
\begingroup
\small
\[
\begin{aligned}
 &(I\times X)
 \mathop{\times}_{\prod_{i=1}^{k}(I_i\times L_i)}
 \left(\prod_{i=1}^{k}(I_i\times P_i)\right)\\
 &=(I_0\times L_0)\times{}^\circ(I\times X)^\circ
 \times\prod_{i=1}^{k}
 \left(I_i\times L_i\times{}^\circ(I_i\times P_i)\right)\\
 &=(I_0\times L_0)\times{}^\circ(I\times X)^\circ
 \times\prod_{i=1}^{k}
 \left(I_i\times L_i\times{}^\circ P_i\right)\\
 &=(-1)^a(I_0\times L_0)
 \times({}^\circ I^\circ\times{}^\circ X^\circ)
 \times\prod_{i=1}^{k}
 \left(I_i\times L_i\times{}^\circ P_i\right)\\
 &=(-1)^{a+b}
 \left(I_0\times{}^\circ I^\circ\times\prod_{i=1}^{k}I_i\right)
 \times L_0\times{}^\circ X^\circ
 \times\prod_{i=1}^{k}(L_i\times{}^\circ P_i)\\
 &=(-1)^{a+b}I\times
 \left(X\mathop{\times}_{\prod_{i=1}^{k}L_i}
 \prod_{i=1}^{k}P_i\right),
\end{aligned}
\]
\endgroup
其中
\[
 b=k\dim X+
 \sum_{i=1}^{k-1}\sum_{j=1}^{i}\dim P_j.
\]
因此
\[
 a+b\equiv
 \sum_{i=1}^{k-1}\sum_{j=1}^{i}(n-\dim P_j)\pmod2.
\]
引理得证。
\end{proof}

下述命题是定义~\ref{def:8-4-1}、定义~\ref{def:8-9-5}
与引理~\ref{lem:8-9-6} 的直接推论。

\begin{proposition}{8.9.7}\label{prop:8-9-7}
作为定向空间，有
\[
\begin{aligned}
 &\mathcal M_{k+1}^{\mathrm{main}}
 (\{J_\rho\}_\rho:\beta;\operatorname{twp}(\rho);
 [0,1]\times P_1,\ldots,[0,1]\times P_k)\\
 &\qquad={}
 [0,1]\times\mathcal M_{k+1}^{\mathrm{main}}
 (J;\beta;P_1,\ldots,P_k).
\end{aligned}
\]
\end{proposition}

\pagebreak
还要注意
\begin{equation}
\begin{aligned}
 &\partial\mathcal M_{k+1}^{\mathrm{main}}
 (\{J_\rho\}_\rho:\beta;\operatorname{twp}(\rho);
 [0,1]\times P_1,\ldots,[0,1]\times P_k)\\
 &\quad={}
 \{1\}\times\mathcal M_{k+1}^{\mathrm{main}}
 (J;\beta;P_1,\ldots,P_k)\\
 &\qquad{}
 -\{0\}\times\mathcal M_{k+1}^{\mathrm{main}}
 (J;\beta;P_1,\ldots,P_k)\\
 &\qquad{}
 -[0,1]\times\partial\mathcal M_{k+1}^{\mathrm{main}}
 (J;\beta;P_1,\ldots,P_k).
\end{aligned}
\tag{8.9.8}\label{eq:8-9-8}
\end{equation}

因此，在
$S([0,1]\times L;\Lambda_{0,\mathrm{nov}})$ 上定义各运算
$\mathfrak m_{k,\beta}$
\translatornote{源书的引句写成“定义 $\mathfrak m_{1,\beta}$”，
但紧随的公式定义 $\mathfrak m_{k,\beta}$，且所得对象是全部高阶运算
组成的滤过 $A_\infty$ 结构；故依公式改正。}如下。对
$(k,\beta)\ne(1,\beta_0)$，置
\[
 \mathfrak m_{k,\beta}(P_1,\ldots,P_k)
 =\Bigl(
 \mathcal M_{k+1}^{\mathrm{main}}
 (\{J_\rho\}_\rho:\beta;\operatorname{twp}(\rho);
 P_1,\ldots,P_k),\operatorname{ev}_0\Bigr),
\]
并置
\[
 \mathfrak m_{1,\beta_0}(P)=(-1)^{n+1}\partial P.
\]
于是，在
$S([0,1]\times L;\Lambda_{0,\mathrm{nov}})$
的某个适当子复形上得到滤过 $A_\infty$ 结构。
（比较定义~\ref{def:3-5-6}。）还可看出，
$\operatorname{Incl}_{1,\beta_0}$ 与这些滤过 $A_\infty$ 结构相容。
严格地说，为定义滤过 $A_\infty$ 运算，必须在逐时积模空间上取扰动。
因此，必须把 $\operatorname{Incl}_{1,\beta_0}$
扩张成有隙滤过 $A_\infty$ 同态 $\operatorname{Incl}$。
